\documentclass[aps,pra,twocolumn,superscriptaddress,longbibliography]{revtex4-2}
\usepackage{amsmath,amssymb,graphicx,xcolor}

\begin{document}

\title{The Geometric Atom: Atomic Structure as a Packing Problem}
\thanks{Reconstructing Atomic Structure from the Principle of Finite Capacity}

\author{J.~Loutey}
\affiliation{Independent Researcher, Kent, Washington}

\date{February 22, 2026}

\begin{abstract}
We present a discrete geometric framework for the hydrogen atom based on finite information capacity. The hydrogen state space is modeled as a paraboloid lattice dual to the SO(4,2) conformal algebra. This single-particle model exhibits a fundamental algebra-geometry duality: \textbf{(i) Algebraic exactness:} \textbf{[PANEL-VERIFIED]} Transition operators $T_\pm, L_\pm$ reproduce the exact Rydberg quantum-number structure ($n$ and $\ell(\ell+1)$) from operator eigenvalues; the energies $E_n = -1/(2n^2)$ follow once the matched kinetic scale is applied (an Observation, Paper~7). \textbf{(ii) Geometric artifacts:} \textbf{[MEASURED]} The graph Laplacian lifts $s/p$ degeneracy at finite lattice cutoff ($37\%$ at $n_{\max}=5$, oscillating through $16\%$ at $8$ and $1.7\%$ at $10$ to $0.39\%$ at $30$ on the binary production lattice; see the \S III QA note) through differential weighted connectivity. Systematic convergence analysis from $n_{\max}=5$ to 30 demonstrates this is a discretization artifact---spectral aliasing from lattice truncation---that oscillates and decays below 1\%, recovering the exact Coulomb degeneracy in the continuum limit. \textbf{(iii) Geometric phase structure:} \textbf{[SYMBOLIC PROOF]} The Berry phase vanishes identically for the real-valued transition operators; the edge-weight log-holonomy $\Theta(n) = -2\ln((n{+}1)/n) \sim n^{-1}$ measures the curvature of the weight function across the lattice. \textbf{Crucially, this model describes only the electron manifold}---electromagnetic interactions and the fine structure constant $\alpha$ cannot emerge from a single-particle lattice and require coupling to photon degrees of freedom, which we address in a companion paper.
\end{abstract}

\maketitle

\section{Introduction: The Texture of the Vacuum}

Physics traditionally assumes spacetime is a smooth continuum. Quantum mechanics inherits this assumption, treating the Hilbert space as infinite-dimensional with continuous operators. Yet this framework produces persistent conceptual difficulties---wave-particle duality, nonlocality, and the measurement problem---that resist resolution within the continuum paradigm.

We propose an alternative foundation: \textit{physical systems encode finite information}. A quantum state labeled by quantum numbers $(n, l, m)$ cannot store infinite positional precision. If information density is finite, then the state space itself must be \textit{discrete}---a lattice, not a continuum. The hydrogen atom becomes a \textbf{graph}: nodes represent quantum states, edges represent physical transitions, and the Hamiltonian becomes a matrix encoding connectivity.

This discrete structure possesses a natural geometry. The hydrogen atom's $SO(4)$ symmetry~\cite{fock1935}, extended to the $SO(4,2)$ dynamical group~\cite{barut1967}, has a discrete dual: the \textbf{paraboloid lattice}, where quantum numbers $(n, l, m)$ map to coordinates on a curved 3D surface. The radius grows parabolically as $r \sim n^2$, and the depth encodes energy: $z = -1/n^2$.

This lattice exhibits a profound duality between \textit{algebra} and \textit{geometry}:
\begin{itemize}
\item \textbf{The Algebraic Skeleton:} Transition operators $T_\pm$ (radial) and $L_\pm$ (angular) generate the spectrum. Their eigenvalues reproduce the exact Rydberg quantum-number structure without correction; the formula $E_n = -1/(2n^2)$ follows once the matched kinetic scale is applied (an Observation, Paper~7).
\item \textbf{The Geometric Flesh:} The lattice \textit{topology}---the pattern of connections---generates physical corrections. High angular momentum states occupy high-connectivity nodes, creating differential weighted connectivity at finite truncation. The edge-weight log-holonomy provides a measure of lattice curvature that decays as $n^{-1}$.
\end{itemize}

In this paper, we present computational evidence for this dual framework. Section~\ref{sec:algebra} demonstrates spectral exactness. Section~\ref{sec:geometry} analyzes the $s/p$ degeneracy breaking by the graph Laplacian, demonstrating it is a discretization artifact that vanishes in the continuum limit. Section~\ref{sec:relativity} analyzes the geometric phase structure, showing that the Berry phase vanishes while the log-holonomy decays as $n^{-1}$. Section~\ref{sec:discussion} discusses the algebra-geometry duality and its connection to wave-particle complementarity. Section~\ref{sec:conclusion} summarizes our findings.


\section{The Algebraic Skeleton: Exact Spectroscopy}
\label{sec:algebra}

The paraboloid lattice is generated by two sets of ladder operators forming the hydrogen atom's dynamical algebra $SU(2) \otimes SU(1,1)$:
\begin{align}
\text{Angular:} \quad & L_z, L_\pm \quad (\text{standard angular momentum}), \\
\text{Radial:} \quad & T_0, T_\pm \quad (\text{energy shell transitions}).
\end{align}
These operators act on basis states $|n, l, m\rangle$ with selection rules:
\begin{align}
L_\pm |n, l, m\rangle &\propto |n, l, m \pm 1\rangle, \\
T_\pm |n, l, m\rangle &\propto |n \pm 1, l, m\rangle.
\end{align}
The operator weights are determined by the Biedenharn-Louck calculus \cite{biedenharn1981} for the $SU(2)$ Clebsch-Gordan coefficients and by the noncompact $SU(1,1)$ ladder structure \cite{barut1967}.

\subsection{Spectral Reproduction}

The critical test of this algebra is spectroscopic: do the operators reproduce the Rydberg formula? We construct the Hamiltonian operator from the Casimir invariants of the algebra. \textbf{[PANEL-VERIFIED]} The energy eigenvalues are
\begin{equation}
E_n = -\frac{1}{2n^2} \quad (\text{Hartree atomic units}),
\label{eq:rydberg}
\end{equation}
in exact agreement with the hydrogen spectrum once the matched kinetic scale $\kappa = -1/16$ fixes the energy unit (matched, not fitted; its coincidence with the Fock Jacobian factor is an Observation, Paper~7). No fitting parameters are required; the algebra determines the quantum-number structure uniquely.

This result validates the lattice as a faithful discrete representation of the hydrogen Hilbert space. The nodes are not a coarse-graining---they \textit{are} the quantum states. The edges encode physically allowed transitions. The quantum-number structure emerges from operator eigenvalues, not from solving differential equations; the energies follow only once the matched scale $\kappa$ is applied (an Observation, Paper~7). (Verified in \texttt{tests/test\_paper1\_rydberg.py}: the $SU(1,1)$ number operator $N = -2[T_+, T_-]$, assembled purely from the ladder generators, has integer spectrum $\{n\}$ and the $SU(2)$ Casimir $L^2$ gives $l(l+1)$, both bit-exact, so $E_n = -1/(2n^2)$ follows from the algebra. This is the dynamical-symmetry construction; the discrete graph Laplacian $L = D - A$ is separately positive-semidefinite and reproduces the spectrum only in the continuum limit, Paper~7~\cite{loutey_paper7}.)


\section{The Geometric Flesh: Topological Artifacts of Finite Truncation}
\label{sec:geometry}

While the algebraic operators yield the exact spectrum, the \textit{geometric structure} of the lattice---its connectivity pattern---introduces systematic artifacts when the lattice is truncated at finite $n_{\max}$. We demonstrate this by constructing the graph Laplacian and analyzing its symmetry-breaking properties as a function of lattice size.

\subsection{The Graph Laplacian Hamiltonian}

In graph theory, the kinetic energy of a particle on a lattice is described by the \textbf{graph Laplacian}~\cite{chung1997}:
\begin{equation}
L = D - A,
\label{eq:laplacian}
\end{equation}
where $A_{ij}$ is the adjacency matrix (edge weights) and $D$ is the weighted degree matrix\footnote{The weighted degree $D_{ii} = \sum_j |A_{ij}|$ represents the total connectivity strength of node $i$. For transition matrices with complex coefficients, we use absolute values to obtain undirected edge weights, explaining the non-integer values in Table~\ref{tab:degrees}.}. This Laplacian is the discrete analog of $-\nabla^2$.

For the paraboloid lattice, we construct $A$ from the transition operators:
\begin{equation}
A = |T_+| + |T_-| + |L_+| + |L_-|,
\end{equation}
where $|\cdot|$ denotes element-wise absolute value (undirected edges). The Hamiltonian becomes
\begin{equation}
H = \beta L + V = \beta (D - A) + V,
\label{eq:ham_graph}
\end{equation}
where $V_{ii} = -1/n_i^2$ is the diagonal Coulomb potential and $\beta$ is a scaling factor.

\subsection{Differential Connectivity by Angular Momentum}

The crucial observation is that the degree matrix $D$ is \textit{not rotationally invariant}. Computing the weighted degrees for states with $n=2$:
\begin{align}
\text{Pole (2s):} \quad & D_{(2,0,0)} = 0.854, \\
\text{Equator (2p):} \quad & D_{(2,1,0)} = 3.416.
\end{align}
\textbf{[PANEL-VERIFIED]} States with higher angular momentum $l$ have \textit{more connections}---more available transitions to neighboring states. This higher connectivity translates to higher on-site energy in the Laplacian formulation.

\subsection{The Splitting as Discretization Artifact}

We performed exact diagonalization of $H$ (Eq.~\ref{eq:ham_graph})
for lattices with $n_{\max}$ ranging from 5 to 30. Using the Lanczos
algorithm, we identified the eigenstates with maximal overlap with
$|2,0,0\rangle$ (2s) and $|2,1,0\rangle$ (2p) and computed the
relative splitting $\Delta E_{\mathrm{rel}} = |E_{2p} - E_{2s}|/|E_{2s}|$.

At $n_{\max} = 10$ (385 nodes), the splitting is:
\begin{align}
\lambda_{2s} &= 0.0202, \\
\lambda_{2p} &= 0.0238, \\
\Delta E &= \lambda_{2p} - \lambda_{2s} = 0.0035,
\end{align}
corresponding to a 16\% relative splitting. However, this degeneracy
breaking is \emph{not} an emergent physical force. Systematic
convergence analysis reveals it is a \textbf{discretization artifact}
of finite lattice truncation---the graph-theoretic analog of spectral
aliasing in finite-element methods.

\paragraph{QA correction (2026-06-14).}  \textbf{[MEASURED]} The specific quantitative
values in this section---the degrees $D_{2s} = 0.854$, $D_{2p} = 3.416$,
the eigenvalues $\lambda_{2s} = 0.0202$, $\lambda_{2p} = 0.0238$, and the
$16\%$ splitting quoted at $n_{\max} = 10$---are \emph{not} reproduced by
the current production lattice (the $A = |T_+| + |T_-| + |L_+| + |L_-|$
CG-magnitude construction they assume is not implemented in
\texttt{geovac/}; a faithful rebuild of that construction gives
$D_{2s} \approx 1.9$ and $D_{2p} \approx 3.8$).  Two lattices must be
kept apart (2026-09-03): on the CG-weighted construction the $2s/2p$
splitting does \emph{not} decay---$129\%$, $68\%$, $84\%$ and $139\%$ at
$n_{\max} = 8, 10, 15, 20$---whereas on the \emph{binary} production
lattice (\texttt{GeometricLattice}, all edge weights 1, the object
\texttt{AtomicSolver} uses) it is $37\%$, $12.7\%$, $0.65\%$, $15.6\%$,
$5.1\%$, $1.7\%$, $2.7\%$, $1.7\%$, $1.1\%$, $4.1\%$, $1.8\%$ and $0.39\%$
at $n_{\max} = 5, 6, 7, 8, 9, 10, 12, 15, 18, 20, 25, 30$.  What \emph{is}
verified (\texttt{tests/test\_trunk\_qa\_splitting.py},
\texttt{test\_trunk\_qa\_degree\_table.py}) is therefore:\ higher-$l$
states have higher connectivity ($D_{2p} > D_{2s}$) on the CG
construction, and on the binary lattice the $s/p$ splitting decays,
strongly oscillating, to below $1\%$ by $n_{\max} = 30$.  The specific
CG-construction figures are retained as illustrative of that
construction only. \emph{This caveat applies wherever these figures
recur---the abstract, Table~I, Appendix~B, and the convergence list
below, which now prints the binary-lattice measurements.  (The
placeholder figure environments that also restated them were removed
2026-09-01.)}

\subsection{Convergence Analysis}

The key diagnostic is the behavior of $\Delta E_{\mathrm{rel}}$ as
$n_{\max} \to \infty$. \textbf{[MEASURED]} We swept the lattice cutoff from $n_{\max} = 5$
(55 vertices) through $n_{\max} = 30$ (9{,}455 vertices) and tracked
the 2s/2p splitting at each resolution:

\begin{itemize}
  \item At $n_{\max} = 5$: $\Delta E_{\mathrm{rel}} \approx 37\%$ (the
    maximum on the grid)
  \item At $n_{\max} = 6, 7$: $12.7\%$, $0.65\%$
  \item At $n_{\max} = 8$: $\approx 16\%$ (a secondary peak); by
    $n_{\max} = 10$ it has fallen to $\approx 1.7\%$
  \item At $n_{\max} = 15, 18$: $1.7\%$, $1.1\%$
  \item At $n_{\max} = 20$: $\Delta E_{\mathrm{rel}} \approx 4.1\%$ (a
    third excursion)
  \item At $n_{\max} = 25, 30$: $1.8\%$, $0.39\%$
\end{itemize}
(binary production lattice, full-spectrum maximum-overlap identification,
\texttt{tests/test\_trunk\_qa\_splitting.py}; the list printed before
2026-09-03---$13\%$, $16\%$, $0.3\%$, $0.005\%$---did not reproduce at
$n_{\max} = 5$, $20$ or $30$.)

The splitting does not decay monotonically; it \emph{oscillates} at
intermediate resolutions before collapsing to negligible values. This
oscillatory decay is the hallmark of spectral aliasing: at low
$n_{\max}$, the truncated boundary of the graph Laplacian introduces
standing-wave reflections that differentially shift eigenvalues in
distinct angular momentum sectors. As the boundary recedes (larger
$n_{\max}$), these reflections weaken and the exact $s/p$ degeneracy
of the continuum Coulomb problem is recovered.

\subsection{Physical Interpretation}

\textbf{[PANEL-VERIFIED]} The 2p state, with its higher weighted degree
($D_{2p} > D_{2s}$; the ``$4 \times$'' ratio came from the
non-reproducing Table~\ref{tab:degrees} figures --- see the \S III QA
note), is more sensitive to the
lattice boundary than the 2s state. At finite $n_{\max}$, the
truncation acts as an absorbing wall that differentially perturbs
high-connectivity nodes. This is a \emph{geometric boundary effect},
not an emergent centrifugal barrier.

Crucially, \textit{no explicit} $l(l+1)/r^2$ term appears in the
Hamiltonian. The transient splitting arises entirely from the finite
extent of the graph, not from any force encoded in the topology.
The graph Laplacian does not spontaneously generate the Lamb shift
or replace electromagnetic interactions; it faithfully reproduces
the Coulomb degeneracy in the limit of large lattices.  \textbf{[MEASURED]} The $s/p$
splitting is a characterisable discretization artifact whose aliasing
boundaries are mapped empirically through $n_{\max} = 30$ (no
convergence rate is proven; see the \S III QA note).


\section{Geometric Phase Structure of the Lattice}
\label{sec:relativity}

The lattice carries a natural notion of parallel transport around closed loops
(plaquettes).  We examine whether this transport produces a nontrivial
geometric phase---a Berry holonomy~\cite{berry1984}---and find that the answer is negative for
the real-valued transition operators used here.  We then characterize the
\emph{log-holonomy}, which measures curvature of the edge-weight function,
and discuss what would be required for a genuine geometric phase.

\subsection{Flat Connection for Real Operators}

A \textbf{rectangular plaquette} on the lattice is constructed via radial
($T_\pm$) and azimuthal ($L_\pm$) operators:
\begin{equation}
|n, l, m\rangle \xrightarrow{T_+} |n+1, l, m\rangle \xrightarrow{L_+}
|n+1, l, m+1\rangle \xrightarrow{T_-} |n, l, m+1\rangle \xrightarrow{L_-}
|n, l, m\rangle.
\label{eq:plaquette}
\end{equation}
The Berry phase for this plaquette is the argument of the holonomy product:
\begin{equation}
\theta = \arg\!\bigl(
  \langle n\!+\!1,l,m|T_+|n,l,m\rangle\,
  \langle n\!+\!1,l,m\!+\!1|L_+|n\!+\!1,l,m\rangle\,
  \langle n,l,m\!+\!1|T_-|n\!+\!1,l,m\!+\!1\rangle\,
  \langle n,l,m|L_-|n,l,m\!+\!1\rangle
\bigr).
\label{eq:berry_arg}
\end{equation}
Each matrix element is a Biedenharn--Louck coupling coefficient.  For the
$SU(2)$ angular-momentum operators,
$\langle l,m\!+\!1|L_+|l,m\rangle = \sqrt{(l-m)(l+m+1)}$,
and for the $SU(1,1)$ radial operators,
$\langle n\!+\!1,l|T_+|n,l\rangle = \sqrt{(n-l)(n+l+1)/4}$.
\textbf{[SYMBOLIC PROOF]} All four factors are \textbf{real and positive}.  Therefore
\begin{equation}
\theta = \arg(\text{positive real number}) = 0
\qquad\text{for every plaquette at every $(n,l,m)$.}
\label{eq:berry_zero}
\end{equation}
The lattice connection defined by these operators is \emph{flat}: it carries
zero Berry curvature.

\textbf{Erratum.}  An earlier version of this paper reported a fitted
exponent $k = 2.113 \pm 0.015$ for the Berry phase scaling.  That value was
a theoretical prediction placed as a placeholder pending numerical
validation; the accompanying figure was never generated (it remained a
\texttt{[Placeholder]}).  Subsequent computation confirms
$\theta = 0$ identically for all plaquettes, as shown above.  The claimed
$k = 2.113$ is therefore \textbf{retracted}.

\subsection{Log-Holonomy of Edge Weights}

Although the Berry phase vanishes, the lattice edge weights carry a
well-defined \emph{log-holonomy}---the logarithmic circulation around a
plaquette:
\begin{equation}
\Theta(n) = \ln w_{12} + \ln w_{23} - \ln w_{34} - \ln w_{41},
\label{eq:log_holonomy}
\end{equation}
where $w_{ij}$ are the edge weights of the \emph{topological-weight}
lattice (\texttt{GeometricLattice(topological\_weights=True)}); on the
binary production lattice and on the CG-magnitude construction of \S III
the $T$-edge weight is $m$-independent and the $L$-edge weight is
$n$-independent, so $\Theta \equiv 0$ there (2026-09-03).  \textbf{[SYMBOLIC PROOF]} For the
topological weights $w = 1/(n_1 n_2)$, the $T_\pm$ edges share the factor
$1/\bigl(n(n+1)\bigr)$ and cancel, leaving
\begin{equation}
\Theta(n) = \ln\!\frac{1}{(n+1)^2} - \ln\!\frac{1}{n^2}
           = -2\ln\!\frac{n+1}{n}
           \;\sim\; -\frac{2}{n}
           \qquad (n \gg 1).
\label{eq:log_holonomy_exact}
\end{equation}
Both results---$\theta = 0$ on every plaquette (with a phased-plaquette
control confirming the detector is not vacuous) and the exact
$\Theta(n)$ closed form with its $n^{-1}$ decay---are pinned by
\texttt{tests/test\_paper1\_geometric\_phase.py} (added 2026-09-01,
closing a coverage gap: the prior Berry-phase artifact was archived
with the retracted $k = 2.113$ claim).
The per-plaquette log-holonomy $\Theta(n) = -2\ln((n+1)/n)$
decays as $n^{-1}$ asymptotically (a closed form, so no fitted exponent or
$R^2$ is involved).  This measures the \textbf{curvature of the weight function}
---how edge weights vary across the lattice---not a genuine geometric phase
in the Berry sense.

\subsection{Toward a Genuine Geometric Phase}

A nonzero Berry phase requires complex-valued transition amplitudes.
Possible extensions include:
\begin{enumerate}
\item the Condon--Shortley phase convention, which attaches $(-1)^m$ to
  $Y_{\ell m}$ and hence flips the sign of alternate $L_\pm$ matrix
  elements (convention-dependent, yielding discrete $\{0,\pi\}$ phases);
\item A $U(1)$ gauge extension---attaching phase factors
  $e^{i f(n,l,m)}$ to the transition operators---which would produce
  continuously varying curvature;
\item An adiabatic Berry connection from eigenstates of a
  parameter-dependent Hamiltonian (as in the hyperspherical $P_{\mu\nu}(R)$
  connection of Ref.~\cite{paper13}).
\end{enumerate}
\textbf{[OPEN]} Whether any such extension recovers a scaling exponent near $k \approx 2$
(the velocity-dependent kinematic factor $v^2 \propto n^{-2}$) remains an
open question for future work.


\section{Discussion: The Algebra-Geometry Duality}
\label{sec:discussion}

Our results reveal a profound duality in the structure of the paraboloid lattice that mirrors the wave-particle duality of quantum mechanics.

\subsection{The Dual Framework}

\begin{enumerate}
\item \textbf{Algebraic Structure (Wave-like):} The operators $T_\pm, L_\pm$ and their eigenvalues reproduce the \textit{unperturbed} hydrogen quantum-number structure exactly; the energies (Eq.~\ref{eq:rydberg}) follow once the matched scale $\kappa$ is applied. This represents the continuous, exact aspect---analogous to the wave description in quantum mechanics.

\item \textbf{Geometric Structure (Particle-like):} The graph Laplacian $L = D - A$ encodes \textit{perturbations}---corrections arising from discrete connectivity. The weighted degree matrix $D$ introduces differential connectivity between angular momentum sectors at finite truncation. The edge-weight log-holonomy decays as $n^{-1}$ (Section~\ref{sec:relativity}), measuring the curvature of the weight function. This represents the discrete aspect of the lattice formulation.
\end{enumerate}

This duality is suggestive: the algebraic operators provide the exact quantum-number structure, while the graph topology introduces finite-size corrections that must be controlled.

\subsection{Why Calibration Fails}

A critical observation: when we attempt to match the graph Laplacian's eigenvalues directly to physical energies, calibration fails. Eigenvalues remain positive (Section~\ref{sec:geometry}), far from the negative binding energies of hydrogen.

This is not a failure---it is a fundamental feature of the duality. The Laplacian encodes \textit{differential} effects (splitting between states), not absolute energies. The algebraic operators provide the reference frame; the geometric operators provide the corrections. Attempting to derive absolute energies from graph topology alone conflates these complementary roles.

\subsection{The Absence of Fine Structure}

\textbf{[MEASURED]} A comprehensive computational search for the fine structure constant $\alpha \approx 1/137$ within the lattice geometry yields a null result. We examined:
\begin{enumerate}
\item \textbf{Operator ratios:} \textbf{[PANEL-VERIFIED]} The spectral-norm gearing ratio
$\|L_+\|/\|T_+\| = 2.0$ \emph{exactly}---at every lattice cutoff
tested, not proved for all $n_{\max}$---while the Frobenius-norm ratio drifts through
$\approx 1.77$ without converging
(\texttt{tests/test\_trunk\_qa\_gearing.py}).  Neither encodes a
physical constant such as $137$ or $1/\alpha$: the spectral value is $2$,
and the drifting one is not a constant.  The commutator
$[T_+, L_+] = 0$, however, \emph{is} bit-exact.
\item \textbf{Commutator defects:} \textbf{[PANEL-VERIFIED]} The operators commute exactly, $[T_+, L_+] = 0$, indicating perfect integrability rather than non-commutative geometry.
\item \textbf{Holonomy defects:} \textbf{[SYMBOLIC PROOF]} The Berry phase is identically zero for real operators (Section~\ref{sec:relativity}); the log-holonomy decays as $n^{-1}$, reflecting the asymptotic flatness of the weight function.
\end{enumerate}

This null result is \textit{physically correct}. The fine structure constant $\alpha = e^2/(4\pi\epsilon_0 \hbar c)$ measures the strength of electromagnetic interactions---specifically, the coupling between electron and photon fields. Our lattice describes the \textit{kinematic} structure of the electron state space without electromagnetic interactions. The hydrogen spectrum $E_n = -1/(2n^2)$ is independent of $\alpha$; fine structure appears only through:
\begin{align}
\text{Relativistic corrections:} \quad & \Delta E_{\text{rel}} \sim \alpha^2 \cdot E_n / n, \\
\text{Spin-orbit coupling:} \quad & \Delta E_{\text{so}} \sim \alpha^2 \cdot E_n \cdot l(l+1) / n^3, \\
\text{Lamb shift (QED):} \quad & \Delta E_{\text{Lamb}} \sim \alpha^5 \cdot mc^2.
\end{align}
These require photon degrees of freedom, which are absent from our single-particle lattice.

This clarifies the role of geometry: the lattice encodes \textit{kinematic} constraints (spectrum, angular momentum structure) but not \textit{dynamical} couplings. To incorporate fine structure, one must introduce a photon lattice and define interaction edges weighted by $\alpha$. The fine structure constant would then appear as the \textit{coupling strength between lattices}, not as an intrinsic geometric property of either lattice alone.

\subsection{Artifact vs.\ Physics: Drawing the Line}

Our convergence analysis establishes a clear separation between
discretization artifacts and genuine physical content of the graph
Laplacian. The $s/p$ splitting is firmly in the artifact category:
it vanishes as $n_{\max} \to \infty$ and exhibits the oscillatory
decay characteristic of finite-size aliasing.

By contrast, the spectral bound $\lambda_{\max} \to 2 d_{\max} = 8$
(hence the extremal eigenvalue $\kappa\lambda_{\max} \to -1/2$; the graph
spectrum is \emph{not} a level-by-level Rydberg ladder,
\texttt{tests/test\_paper7\_graph\_convergence.py}) and the edge-weight log-holonomy
($\Theta \propto n^{-1}$, Section~\ref{sec:relativity}) are
\emph{convergent} properties that survive the continuum limit.
These represent genuine structural content encoded in the topology.

This distinction is essential: the graph Laplacian does not generate
electromagnetic forces or the Lamb shift. It encodes the
\textit{kinematic} structure of the electron state space. Finite
lattice effects must be understood as numerical artifacts, not
new physics.

\subsection{The Wavefunction as Distribution}

If quantum states are discrete nodes, what is the wavefunction $\psi(r)$? One possible interpretation is that continuous wavefunctions correspond to statistical distributions over discrete lattice sites---a coarse-graining in which the Schr\"odinger equation becomes a diffusion equation on the graph. Under this reading, photon absorption would transfer discrete angular momentum via edge traversal. However, this interpretation is not required by the mathematical framework presented here; the graph-theoretic formulation is compatible with standard quantum-mechanical interpretation of $\psi$ as a probability amplitude.


\section{Conclusion}
\label{sec:conclusion}

We have demonstrated that quantum mechanics, when formulated on a discrete paraboloid lattice, exhibits a dual structure:
\begin{itemize}
\item \textbf{Algebraic exactness:} \textbf{[PANEL-VERIFIED]} Transition operators reproduce the hydrogen quantum-number structure without corrections; the energies follow with the matched kinetic scale (Paper~7).
\item \textbf{Understood discretization artifacts:} \textbf{[MEASURED]} The graph Laplacian lifts $s/p$ degeneracy at finite $n_{\max}$ (up to $37\%$ at $n_{\max}=5$ on the binary production lattice, $16\%$ at $8$, $1.7\%$ at $10$), but systematic convergence analysis confirms this is spectral aliasing from lattice truncation---not an emergent physical force. The splitting oscillates ($4.1\%$ at $n_{\max}=20$) and decays below 1\% by $n_{\max}=30$ ($0.39\%$), recovering the exact Coulomb degeneracy; on the CG-weighted construction of \S III the lift does \emph{not} decay (measured 2026-09-03), so the artifact-decay claim is scoped to the binary lattice.
\item \textbf{Geometric phase structure:} \textbf{[SYMBOLIC PROOF]} The Berry phase vanishes identically for the real-valued $SU(2)/SU(1,1)$ operators; the edge-weight log-holonomy decays as $\Theta \propto n^{-1}$, measuring the curvature of the weight function. Whether a $U(1)$ gauge extension can produce a genuine geometric phase with relativistic scaling remains an open question.
\item \textbf{Kinematic completeness:} \textbf{[OBSERVATION]} The lattice encodes the $\alpha$-independent single-electron structure---the Coulomb spectrum, angular momentum structure, and geometric phase effects---while cleanly separating boundary artifacts from physical content.
\end{itemize}

These results support a graph-theoretic reformulation of the hydrogen state space in which the algebraic operators encode the exact quantum-number structure (the energies following with the matched kinetic scale) and the graph Laplacian introduces controlled, quantitatively understood finite-size artifacts. The algebra-geometry duality is suggestive of a connection to wave-particle complementarity, though we do not claim it establishes one. Critically, the graph Laplacian does not replace electromagnetic interactions or generate the Lamb shift; it encodes the kinematic manifold of the electron.


\section{Limitations and Next Steps: Missing Photon Degrees of Freedom}
\label{sec:limitations}

\textbf{What This Model Does Not Include:}

The paraboloid lattice presented here is explicitly a \textit{single-particle} model describing only the electron's quantum state manifold. It does not include:
\begin{enumerate}
\item \textbf{Photon degrees of freedom:} The electromagnetic field ($U(1)$ gauge symmetry) is absent.
\item \textbf{Electron-photon coupling:} There is no mechanism for radiative transitions or electromagnetic interactions.
\item \textbf{The fine structure constant $\alpha$:} As a pure electron-only model, $\alpha$ cannot appear. The fine structure constant measures the strength of electron-photon coupling---it describes the relationship between \textit{two} manifolds, not a property of either one alone.
\end{enumerate}

\textbf{What Must Be Added:}

To incorporate electromagnetic interactions, one must construct a second geometric structure:
\begin{itemize}
\item \textbf{A photon gauge fiber:} The $U(1)$ electromagnetic gauge symmetry suggests a circular or helical fiber structure attached to each electron state $(n,l,m)$.
\item \textbf{Coupling rules:} Edges connecting the electron lattice to the photon fiber must be defined. The weight of these inter-manifold connections would encode the coupling strength.
\item \textbf{Symplectic impedance:} If both manifolds carry action (units of $\hbar$), their ratio may define a dimensionless coupling constant.
\end{itemize}

\textbf{[CONJECTURE]} We conjecture that the fine structure constant $\alpha$ emerges as the \textit{geometric impedance}---the ratio of action densities between the electron paraboloid and the photon fiber. \textbf{[OBSERVATION]} A companion paper \cite{companion_alpha} records the related
observation that $\alpha$ satisfies the cubic
$\alpha^3 - K\alpha + 1 = 0$ with $K = \pi(B + F - \Delta)$ on the
Hopf fibration.  It concludes that $\alpha$ behaves as a projection
constant, not derivable from any single principle; the impedance
picture above thus remains unexplored.

\textbf{Outlook:}

This framework invites generalization. If hydrogen is a paraboloid, what geometries describe molecules, nuclei, or quantum fields? The answer may lie in higher-dimensional lattices---discretizations of larger symmetry groups.


\begin{acknowledgments}
We thank the developers of \texttt{scipy.sparse} for enabling efficient eigenvalue computations on large lattice Hamiltonians. We acknowledge the foundational work on hydrogen dynamical symmetry by Fock, Barut, and the Biedenharn-Louck formulation of angular momentum coupling.
\end{acknowledgments}

\appendix

\section{Mathematical Definitions}
\label{app:berry}

\subsection{Weighted Degree Matrix}

The weighted degree matrix $D$ is a diagonal matrix where each diagonal element represents the total connection strength of a node:
\begin{equation}
D_{ii} = \sum_{j} |A_{ij}|,
\end{equation}
where $A_{ij}$ is the adjacency matrix constructed from transition operators. For complex-valued operators, the absolute value ensures symmetric (undirected) edge weights. This definition explains why degree values are non-integer (Table~\ref{tab:degrees})---they represent weighted sums of transition matrix elements, not simple edge counts.

\subsection{Berry Phase and Log-Holonomy}

For a plaquette (Eq.~\ref{eq:plaquette}), the Berry phase is
$\theta = \arg(\prod_i \langle s_i | \hat{U}_i | s_{i+1}\rangle)$.
As shown in Section~\ref{sec:relativity}, all four matrix elements are real
and positive, so $\theta = 0$ identically.

The log-holonomy (Eq.~\ref{eq:log_holonomy}) provides a nonzero measure of
weight-function curvature.  For topological weights $w = 1/(n_1 n_2)$, the
exact result is $\Theta(n) = -2\ln\bigl((n+1)/n\bigr)$, decaying as
$n^{-1}$.


\section{Computational Details}

\subsection{Lattice Construction}

The paraboloid lattice was constructed for $1 \leq n \leq 10$, yielding 385 quantum states $|n, l, m\rangle$ with $0 \leq l < n$ and $-l \leq m \leq l$. Transition operators $T_\pm$ and $L_\pm$ were built using Biedenharn-Louck coupling coefficients for $SU(2)$ and $SU(1,1)$ representations.

\subsection{Graph Laplacian Diagonalization}

The adjacency matrix $A$ was computed as the sum of absolute values of all transition operators:
\begin{equation}
A = |T_+| + |T_-| + |L_+| + |L_-|.
\end{equation}
The graph Laplacian $L = D - A$ was stored in compressed sparse row (CSR) format. The weighted degree matrix $D$ of the CG-magnitude construction ranges from $0.71$ to $19.9$ with mean degree $12.0$ at $n_{\max} = 10$ (faithful rebuild, 2026-09-03; the figures printed earlier, $0.854$--$20.624$ and $12.42$, belong to the non-reproducing table of \S III).

Eigenvalue computation used the shift-invert Lanczos method (\texttt{scipy.sparse.linalg.eigsh}) with $\sigma = -1.0$, computing the lowest 20 eigenvalues. Diagonalization required 0.15 seconds on a standard workstation; \textbf{[MEASURED]} the sparse-eigenvalue timing is guarded sub-quadratic and the adjacency nonzero count near-linear (measured exponent $1.07$ over $n_{\max} = 5$--$30$) by \texttt{tests/test\_ov\_scaling\_rigorous.py} (multi-point log--log fits; the frozen bounds are $< 1.8$ on timing and $< 1.15$ on the nonzero count).

State identification was performed by computing overlap probabilities $|\langle n,l,m | \psi_k \rangle|^2$ for each eigenvector $|\psi_k\rangle$, selecting the eigenvector with maximal overlap for each target state.

\subsection{Log-Holonomy Calculation}

Rectangular plaquettes (Eq.~\ref{eq:plaquette}) were identified by searching
for closed paths in $(n,m)$ space at fixed $l$.  The Berry phase
$\theta = \arg(\cdot)$ is identically zero for all plaquettes because every
$SU(2)/SU(1,1)$ matrix element is real and positive
(Section~\ref{sec:relativity}).  The log-holonomy
$\Theta(n) = -2\ln\bigl((n+1)/n\bigr)$ was computed analytically and
verified numerically; its closed form decays as $n^{-1}$ asymptotically (no fitted exponent is involved---a finite-$n$ log--log fit gives $-0.99$ on $n = 10$--$320$).


\section{Supplementary Data}

\begin{table}[h]
\centering
\caption{Weighted node degrees for select quantum states, demonstrating differential connectivity between $s$ and $p$ orbitals. The non-integer values arise from weighted sums of transition matrix elements (see Appendix~\ref{app:berry}).}
\begin{tabular}{cccc}
\hline
State $(n,l,m)$ & Weighted Degree $D_{ii}$ & Type & Energy (Laplacian) \\
\hline
$(1,0,0)$ & 0.854 & 1s & (ground state) \\
$(2,0,0)$ & 0.854 & 2s & $0.0202$ \\
$(2,1,0)$ & 3.416 & 2p & $0.0238$ \\
$(3,0,0)$ & 0.854 & 3s & --- \\
$(3,1,0)$ & 4.270 & 3p & --- \\
$(3,2,0)$ & 6.124 & 3d & --- \\
\hline
\end{tabular}
\label{tab:degrees}
\end{table}

% Three [Placeholder] figure environments (fig:lattice, fig:killswitch,
% fig:relativity) were removed 2026-09-01 (trunk-FULL remediation B10):
% never generated, unreferenced in text, and this paper's own Erratum
% retracted a claim precisely because its figure remained a placeholder.


\begin{thebibliography}{99}

\bibitem{barut1967}
A.~O. Barut and H. Kleinert,
``Transition probabilities of the hydrogen atom from noncompact dynamical groups,''
Phys. Rev. \textbf{156}, 1541 (1967).

\bibitem{fock1935}
V. Fock,
``Zur Theorie des Wasserstoffatoms,''
Z. Phys. \textbf{98}, 145 (1935).

\bibitem{biedenharn1981}
L.~C. Biedenharn and J.~D. Louck,
\textit{Angular Momentum in Quantum Physics},
Encyclopedia of Mathematics and its Applications, Vol. 8 (Addison-Wesley, Reading, MA, 1981).

\bibitem{berry1984}
M.~V. Berry,
``Quantal phase factors accompanying adiabatic changes,''
Proc. R. Soc. Lond. A \textbf{392}, 45 (1984).

\bibitem{chung1997}
F.~R.~K. Chung,
\textit{Spectral Graph Theory},
CBMS Regional Conference Series in Mathematics, Vol. 92 (American Mathematical Society, Providence, RI, 1997).

\bibitem{companion_alpha}
J. Loutey,
``The Fine Structure Constant as a Spectral Coincidence on the Hopf Fibration,''
(companion paper, 2026).

\bibitem{paper13}
J. Loutey,
``The Hyperspherical Lattice: Two-Electron Atoms as Coupled Channel Graphs,''
(companion paper, 2026).

\bibitem{loutey_paper7}
J.~Loutey,
``The Dimensionless Vacuum: Recovering the Schr\"odinger Equation from Scale-Invariant Graph Topology,''
GeoVac Paper~7 (2026).
\end{thebibliography}

\end{document}
